\documentclass[11pt]{article}

\usepackage[margin=1in]{geometry}
\usepackage{amsmath,amssymb,amsthm,mathtools}
\usepackage{bm}
\usepackage{hyperref}

\newtheorem{proposition}{Proposition}
\newtheorem{assumption}{Assumption}
\newtheorem{remark}{Remark}

\newcommand{\diag}{\operatorname{diag}}
\newcommand{\cov}{\operatorname{cov}}
\newcommand{\var}{\operatorname{var}}
\newcommand{\E}{\mathbb{E}}

\title{Precision-Space Mini-Batch Updates for TAGI}
\author{}
\date{}

\begin{document}
\maketitle

\paragraph{Motivation}
In standard Tractable Approximate Gaussian Inference (TAGI)
\cite{goulet2021tractable}, the observation model is
\begin{equation}
    \bm{y} = \bm{z}^{(\mathtt{O})} + \bm{v},
    \qquad
    \bm{V} \sim \mathcal{N}(\bm{0},\bm{\Sigma}_{\bm{V}}).
    \label{eq:obs-model}
\end{equation}
The forward pass gives
\begin{equation}
    \bm{\mu}_{\bm{Y}} = \bm{\mu}_{\bm{Z}^{(\mathtt{O})}},
    \qquad
    \bm{\Sigma}_{\bm{Y}}
    =
    \bm{\Sigma}_{\bm{Z}^{(\mathtt{O})}} + \bm{\Sigma}_{\bm{V}}.
    \label{eq:predictive-y}
\end{equation}
For diagonal output covariance, the innovation quantities used in the
implementation can be written element-wise as
\begin{equation}
    \delta_{\mu,i}
    =
    \frac{y_i-\mu_{Y_i}}{\sigma^2_{Y_i}},
    \qquad
    \delta_{\sigma,i}
    =
    -\frac{1}{\sigma^2_{Y_i}},
    \label{eq:output-innovation}
\end{equation}
where
\(
\sigma^2_{Y_i}
=
\sigma^2_{Z^{(\mathtt{O})}_i}+\sigma^2_{V_i}.
\)
These quantities are propagated backward through the network using the
layer-wise RTS equations. For a scalar parameter
\(
\theta_p^{(j)}
\)
with prior
\(
\theta_p^{(j)} \sim
\mathcal{N}(\mu_{\theta_p},\sigma^2_{\theta_p})
\),
the usual additive TAGI update can be written as
\begin{equation}
    \mu_{\theta_p|\bm{y}}
    =
    \mu_{\theta_p}
    +
    \Delta \mu_{\theta_p}^{\mathrm{add}},
    \qquad
    \sigma^2_{\theta_p|\bm{y}}
    =
    \sigma^2_{\theta_p}
    +
    \Delta \sigma_{\theta_p}^{2,\mathrm{add}}.
    \label{eq:additive-tagi-update}
\end{equation}
Under the diagonal covariance approximation, the accumulated increments have
the local form
\begin{equation}
    \Delta \mu_{\theta_p}^{\mathrm{add}}
    =
    \sigma^2_{\theta_p}\eta_{\theta_p},
    \qquad
    \Delta \sigma_{\theta_p}^{2,\mathrm{add}}
    =
    -\sigma^4_{\theta_p}\Lambda_{\theta_p},
    \label{eq:additive-natural-form}
\end{equation}
where
\(
\eta_{\theta_p}
\)
is a first-order information vector contribution and
\(
\Lambda_{\theta_p}\geq 0
\)
is a scalar precision contribution. For a mini-batch, TAGI implementations
usually obtain
\(
\eta_{\theta_p}
\)
and
\(
\Lambda_{\theta_p}
\)
by summing contributions over the examples in the batch while keeping the
prior moments fixed during the backward pass.

\paragraph{The Mini-Batch Difficulty}
For a single observation, the additive variance formula in
\eqref{eq:additive-natural-form} is the first-order Gaussian conditioning
update used by TAGI. For a mini-batch, however, all examples are processed
with the same frozen prior variance
\(
\sigma^2_{\theta_p}
\).
If the summed information
\(
\sigma^2_{\theta_p}\Lambda_{\theta_p}
\)
is large, the additive update
\[
    \sigma^2_{\theta_p}
    -
    \sigma^4_{\theta_p}\Lambda_{\theta_p}
\]
can become too small or even non-positive. In practice this motivates
heuristics such as a cap factor, which limits the magnitude of the parameter
increments. The purpose of the precision-space update is to replace this cap
with the exact form of the local Gaussian batch posterior.

\paragraph{Local Linear-Gaussian Batch Model}
Consider one scalar parameter
\(
\theta_p
\equiv \theta_p^{(j)}
\)
inside the network and assume that, conditionally on all other random
variables being represented by their current Gaussian moments, each scalar
observation in a mini-batch admits the local model
\begin{equation}
    y_{b,i}
    =
    \mu_{Y_{b,i}}
    +
    h_{b,i,p}(\theta_p-\mu_{\theta_p})
    +
    \varepsilon_{b,i},
    \qquad
    \varepsilon_{b,i}
    \sim
    \mathcal{N}(0,r_{b,i}),
    \label{eq:local-linear-model}
\end{equation}
where
\(
b\in\mathcal{B}
\)
indexes the mini-batch and
\(
i
\)
indexes the output dimension or selected observation node. The quantity
\(
h_{b,i,p}
\)
is the local sensitivity of the predictive mean with respect to
\(
\theta_p
\),
and
\(
r_{b,i}
\)
is the local likelihood variance assigned to that observation. In the usual
TAGI diagonal implementation this role is played by the predictive scalar
variance
\(
\sigma^2_{Y_{b,i}}
\)
from \eqref{eq:predictive-y}. Define
\begin{equation}
    e_{b,i}
    =
    y_{b,i}-\mu_{Y_{b,i}},
    \qquad
    \eta_{\theta_p}
    =
    \sum_{b\in\mathcal{B}}\sum_i
    \frac{h_{b,i,p}e_{b,i}}{r_{b,i}},
    \qquad
    \Lambda_{\theta_p}
    =
    \sum_{b\in\mathcal{B}}\sum_i
    \frac{h_{b,i,p}^2}{r_{b,i}}.
    \label{eq:eta-lambda}
\end{equation}

\begin{proposition}[Exact local batch posterior]
Under the local model \eqref{eq:local-linear-model} and the scalar Gaussian
prior
\(
\theta_p\sim
\mathcal{N}(\mu_{\theta_p},\sigma^2_{\theta_p})
\),
the posterior marginal for
\(
\theta_p
\)
is Gaussian with
\begin{equation}
    \sigma^2_{\theta_p|\mathcal{B}}
    =
    \frac{\sigma^2_{\theta_p}}
    {1+\sigma^2_{\theta_p}\Lambda_{\theta_p}},
    \qquad
    \mu_{\theta_p|\mathcal{B}}
    =
    \mu_{\theta_p}
    +
    \frac{\sigma^2_{\theta_p}\eta_{\theta_p}}
    {1+\sigma^2_{\theta_p}\Lambda_{\theta_p}}.
    \label{eq:exact-local-posterior}
\end{equation}
\end{proposition}

\begin{proof}
Let
\(
\tilde{\theta}_p=\theta_p-\mu_{\theta_p}
\).
The log posterior is, up to an additive constant,
\begin{align}
    \log p(\theta_p|\mathcal{B})
    &=
    -\frac{1}{2}
    \frac{\tilde{\theta}_p^2}{\sigma^2_{\theta_p}}
    -
    \frac{1}{2}
    \sum_{b\in\mathcal{B}}\sum_i
    \frac{(e_{b,i}-h_{b,i,p}\tilde{\theta}_p)^2}{r_{b,i}}  \nonumber \\
    &=
    -\frac{1}{2}
    \left(
    \frac{1}{\sigma^2_{\theta_p}}
    +
    \Lambda_{\theta_p}
    \right)
    \tilde{\theta}_p^2
    +
    \eta_{\theta_p}\tilde{\theta}_p
    + \mathrm{const}.
    \label{eq:posterior-complete-square}
\end{align}
Completing the square gives posterior precision
\(
(\sigma^2_{\theta_p})^{-1}+\Lambda_{\theta_p}
\)
and posterior mean displacement
\(
\eta_{\theta_p}/((\sigma^2_{\theta_p})^{-1}+\Lambda_{\theta_p})
\),
which is exactly \eqref{eq:exact-local-posterior}.
\end{proof}

\paragraph{Precision-Space TAGI Update}
Equation \eqref{eq:exact-local-posterior} suggests replacing the capped
additive mini-batch update with
\begin{equation}
    d_{\theta_p}
    =
    1+\sigma^2_{\theta_p}\Lambda_{\theta_p}.
    \label{eq:precision-denom}
\end{equation}
Using the TAGI additive increments in \eqref{eq:additive-natural-form}, this
can be written without explicitly forming
\(
\eta_{\theta_p}
\)
or
\(
\Lambda_{\theta_p}
\):
\begin{equation}
    d_{\theta_p}
    =
    1
    -
    \frac{
    \Delta \sigma_{\theta_p}^{2,\mathrm{add}}
    }{
    \sigma^2_{\theta_p}
    }.
    \label{eq:denom-from-delta}
\end{equation}
The proposed mini-batch update is therefore
\begin{equation}
\boxed{
\begin{aligned}
    \mu_{\theta_p|\mathcal{B}}
    &=
    \mu_{\theta_p}
    +
    \frac{
    \Delta \mu_{\theta_p}^{\mathrm{add}}
    }{
    d_{\theta_p}
    },
    \\
    \sigma^2_{\theta_p|\mathcal{B}}
    &=
    \frac{
    \sigma^2_{\theta_p}
    }{
    d_{\theta_p}
    },
    \qquad
    d_{\theta_p}
    =
    1
    -
    \frac{
    \Delta \sigma_{\theta_p}^{2,\mathrm{add}}
    }{
    \sigma^2_{\theta_p}
    }.
\end{aligned}}
\label{eq:precision-tagi-update}
\end{equation}
This is the precision-space analogue of the usual TAGI parameter update. It
uses the same forward pass, the same observation innovation
\eqref{eq:output-innovation}, and the same RTS backward pass. Only the final
assimilation of the accumulated parameter increments is changed.

\paragraph{Optional Tempered Form}
If one wants to explicitly temper the amount of mini-batch information, one
may introduce a scalar
\(
\rho\in(0,1]
\)
and use
\begin{equation}
\begin{aligned}
    d_{\theta_p}(\rho)
    &=
    1
    -
    \rho
    \frac{
    \Delta \sigma_{\theta_p}^{2,\mathrm{add}}
    }{
    \sigma^2_{\theta_p}
    },
    \\
    \mu_{\theta_p|\mathcal{B}}
    &=
    \mu_{\theta_p}
    +
    \rho
    \frac{
    \Delta \mu_{\theta_p}^{\mathrm{add}}
    }{
    d_{\theta_p}(\rho)
    },
    \\
    \sigma^2_{\theta_p|\mathcal{B}}
    &=
    \frac{\sigma^2_{\theta_p}}{d_{\theta_p}(\rho)}.
\end{aligned}
\label{eq:tempered-precision-update}
\end{equation}
The default choice is
\(
\rho=1
\),
which introduces no additional hyperparameter and corresponds to full local
Bayesian assimilation of the mini-batch. Values
\(
\rho<1
\)
can be interpreted as fractional likelihood or power posterior updates, but
they are not required by the proposed method.

\paragraph{Consistency With Standard TAGI}
\begin{proposition}[First-order agreement]
For small local information
\(
\sigma^2_{\theta_p}\Lambda_{\theta_p}\ll 1
\),
the precision-space update \eqref{eq:precision-tagi-update} agrees with the
standard additive TAGI update to first order.
\end{proposition}

\begin{proof}
Using
\(
(1+x)^{-1}=1-x+\mathcal{O}(x^2)
\)
with
\(
x=\sigma^2_{\theta_p}\Lambda_{\theta_p}
\),
\begin{align}
    \sigma^2_{\theta_p|\mathcal{B}}
    &=
    \frac{\sigma^2_{\theta_p}}{1+x}
     =
    \sigma^2_{\theta_p}
    -
    \sigma^4_{\theta_p}\Lambda_{\theta_p}
    +
    \mathcal{O}(x^2),
    \\
    \mu_{\theta_p|\mathcal{B}}
    &=
    \mu_{\theta_p}
    +
    \frac{\sigma^2_{\theta_p}\eta_{\theta_p}}{1+x}
     =
    \mu_{\theta_p}
    +
    \sigma^2_{\theta_p}\eta_{\theta_p}
    +
    \mathcal{O}(x\,\sigma^2_{\theta_p}\eta_{\theta_p}).
\end{align}
The leading terms are exactly the additive TAGI increments in
\eqref{eq:additive-natural-form}.
\end{proof}

\paragraph{Positivity and Stability}
\begin{proposition}[Positive variance]
If
\(
\sigma^2_{\theta_p}>0
\)
and
\(
\Lambda_{\theta_p}\geq 0
\),
then the precision-space update satisfies
\begin{equation}
    0
    <
    \sigma^2_{\theta_p|\mathcal{B}}
    \leq
    \sigma^2_{\theta_p}.
    \label{eq:positive-monotone-var}
\end{equation}
\end{proposition}

\begin{proof}
Since
\(
\Lambda_{\theta_p}\geq 0
\),
the denominator
\(
1+\sigma^2_{\theta_p}\Lambda_{\theta_p}
\)
is at least one. Dividing the positive prior variance by this denominator
gives a positive posterior variance that cannot exceed the prior variance.
\end{proof}

This property is the central practical advantage over the additive mini-batch
variance update. The additive update can become non-positive when
\(
\sigma^2_{\theta_p}\Lambda_{\theta_p}>1
\).
The precision update remains positive for any finite non-negative
\(
\Lambda_{\theta_p}
\).
Moreover, the mean update is automatically damped by the same denominator
that increases posterior precision:
\begin{equation}
    \Delta \mu_{\theta_p}^{\mathrm{prec}}
    =
    \frac{
    \sigma^2_{\theta_p}\eta_{\theta_p}
    }{
    1+\sigma^2_{\theta_p}\Lambda_{\theta_p}
    }.
    \label{eq:self-normalized-mean}
\end{equation}
Thus, when the mini-batch contains a large amount of information about a
parameter, the method reduces both its variance and its mean displacement in
a coupled Bayesian way rather than clipping the two quantities independently.

\paragraph{Relation to the Cap Factor}
The capped TAGI update can be written schematically as
\begin{equation}
\begin{aligned}
    \mu_{\theta_p}^{+}
    &=
    \mu_{\theta_p}
    +
    \operatorname{sign}(\Delta\mu_{\theta_p}^{\mathrm{add}})
    \min\left(
    |\Delta\mu_{\theta_p}^{\mathrm{add}}|,
    \frac{\sigma_{\theta_p}}{c_B}
    \right),
    \\
    (\sigma^2_{\theta_p})^{+}
    &=
    \sigma^2_{\theta_p}
    +
    \operatorname{sign}(\Delta\sigma_{\theta_p}^{2,\mathrm{add}})
    \min\left(
    |\Delta\sigma_{\theta_p}^{2,\mathrm{add}}|,
    \frac{\sigma_{\theta_p}}{c_B}
    \right),
\end{aligned}
\label{eq:capped-update}
\end{equation}
where
\(
c_B
\)
is a batch-size-dependent cap factor. This stabilizes training, but it is not
the posterior of a Gaussian conditioning problem. In contrast,
\eqref{eq:precision-tagi-update} is the exact posterior of the local
linear-Gaussian batch model in Proposition 1. The cap factor is therefore
replaced by the posterior precision denominator
\(
d_{\theta_p}
\),
which is determined by the inferred information
\(
\Lambda_{\theta_p}
\)
rather than by a manually selected batch-size table.

\paragraph{Correctness Statement}
The proposed method is correct in the following precise sense. Under the same
diagonal covariance and local moment-matching assumptions used by TAGI, and
under the additional local linear-Gaussian model
\eqref{eq:local-linear-model} for each scalar parameter marginal, the update
\eqref{eq:precision-tagi-update} is the exact Gaussian posterior marginal for
that scalar parameter after assimilating the mini-batch. When the local
information is small, it reduces to the standard additive TAGI update to first
order. For arbitrary finite mini-batch information, it preserves positivity of
all parameter variances by construction.

\paragraph{Viability in Deep Networks}
The update is viable for deep TAGI networks for the same reasons that the
standard layer-wise TAGI update is viable:
\begin{enumerate}
    \item The forward pass still uses the Gaussian Multiplicative
    Approximation and analytic activation moment propagation.
    \item The backward pass still uses the RTS layer-wise conditional update
    and the diagonal covariance approximation.
    \item The computational cost remains linear in the number of learnable
    scalar parameters because \eqref{eq:precision-tagi-update} is element-wise.
    \item The variance update cannot become negative, so no variance collapse
    through additive over-subtraction is possible.
    \item The mean update is self-normalized by the inferred batch precision,
    reducing the need for an external cap factor.
\end{enumerate}
The main approximation is that the sensitivities
\(
h_{b,i,p}
\)
and predictive variances
\(
r_{b,i}
\)
are held fixed while the whole mini-batch is assimilated. This is the same
frozen-prior approximation already present in mini-batch TAGI. The
precision-space update improves this approximation by using the batch
posterior precision implied by the frozen local model, rather than applying a
first-order variance update outside its stable regime.

\paragraph{Role of the Observation Noise}
The observation covariance
\(
\bm{\Sigma}_{\bm{V}}
\)
remains part of the likelihood through
\(
\bm{\Sigma}_{\bm{Y}}
\)
in \eqref{eq:predictive-y}. A smaller
\(
\bm{\Sigma}_{\bm{V}}
\)
increases both the innovation magnitude and the inferred precision
\(
\Lambda_{\theta_p}
\);
a larger
\(
\bm{\Sigma}_{\bm{V}}
\)
decreases both. The precision-space update therefore does not remove the need
for a meaningful observation model. Instead, it separates two roles that are
often conflated in practice:
\begin{equation}
    \bm{\Sigma}_{\bm{V}}
    \quad\text{controls likelihood noise,}
    \qquad
    d_{\theta_p}
    \quad\text{controls Bayesian batch assimilation.}
\end{equation}
This separation is desirable: the observation noise should represent
aleatory uncertainty or model mismatch, whereas the amount of parameter
movement should follow from posterior precision.

\paragraph{Summary}
The precision-space mini-batch TAGI update is
\[
    \mu_{\theta_p|\mathcal{B}}
    =
    \mu_{\theta_p}
    +
    \frac{
    \Delta \mu_{\theta_p}^{\mathrm{add}}
    }{
    1
    -
    \Delta \sigma_{\theta_p}^{2,\mathrm{add}}/
    \sigma^2_{\theta_p}
    },
    \qquad
    \sigma^2_{\theta_p|\mathcal{B}}
    =
    \frac{
    \sigma^2_{\theta_p}
    }{
    1
    -
    \Delta \sigma_{\theta_p}^{2,\mathrm{add}}/
    \sigma^2_{\theta_p}
    }.
\]
It is exact for the local linear-Gaussian batch approximation, first-order
consistent with standard TAGI, positive by construction, and computationally
linear in the number of parameters. It replaces the cap factor by a Bayesian
precision denominator determined by the mini-batch information itself.

\end{document}
